\section{Algorithms and the enumeration of graph classes}\label{sec:algorithms}

The table requires decisions about a graph and about the neighbourhoods
that occur during its reductions. We first give exact procedures for
these decisions. We then describe how to enumerate isomorphism classes
and how to reduce the classification at order $14$ to a smaller family.

\subsection{Cliques and homology}

Fix an ordering of the vertices of a nonempty graph $G$. Starting with
each vertex, extend an increasing clique by adjoining a larger vertex
adjacent to all its members, recording each clique as it is formed.
Continue until no extension is possible.
Every clique has a unique increasing order and a unique sequence of
initial segments, so this procedure lists it exactly once. In particular,
it determines all dimensions of $\cl(G)$, not only its triangles.

For an increasing clique $[v_0,\ldots,v_i]$, use the boundary
\[
 \partial_i[v_0,\ldots,v_i]
   =\sum_{j=0}^i(-1)^j[v_0,\ldots,\widehat v_j,\ldots,v_i].
\]
Adjoin the augmentation $\partial_0:C_0\to\ZZ$, which sends every
vertex to $1$. Write $f_i=\operatorname{rank} C_i$ and let $r_i(R)$
denote the rank of $\partial_i$ over the field $R$. Since consecutive
boundaries compose to zero, elementary linear algebra gives
\begin{equation}\label{eq:betti}
 \dim_R\widetilde H_i(G;R)=f_i-r_i(R)-r_{i+1}(R)
 \qquad(i\ge0),
\end{equation}
where $r_0(R)=1$ and the rank above the highest dimension is zero.
Thus elimination over $\mathbb F_2$ decides membership in
$\mathcal A_{\mathbb F_2}$.
The equality $\chi(G)=1$ is a necessary preliminary condition, but is
not a substitute for the homology calculation.

Integral homology also requires torsion information. If $d_{i+1,j}$
are the nonzero Smith factors of $\partial_{i+1}$, then
\begin{equation}\label{eq:integral-homology}
 \widetilde H_i(G;\ZZ)\cong
 \ZZ^{\,f_i-r_i(\mathbb Q)-r_{i+1}(\mathbb Q)}
 \oplus\bigoplus_{d_{i+1,j}>1}\ZZ/d_{i+1,j}.
\end{equation}
Indeed, the augmented chain complex gives a short exact sequence
\[
 0\longrightarrow\widetilde H_i(G;\ZZ)
 \longrightarrow C_i/\operatorname{im}\partial_{i+1}
 \longrightarrow\operatorname{im}\partial_i\longrightarrow0.
\]
Its last term is free abelian, so the sequence splits. The Smith
factors describe the middle term, and the two adjacent rational ranks
give the displayed free rank. Formula~\eqref{eq:integral-homology},
in every dimension present, therefore decides integral acyclicity.

\subsection{Strong and vertex reductions}

Let $S(G)$ be a truth value. Set $S(\emptyset)=\mathrm{false}$ and
$S(K_1)=\mathrm{true}$. For a graph with at least two vertices evaluate
\begin{align}
 S(G)={}&\bigvee_{v\in V(G)}
 \bigl(S(\Nb_G(v))\land S(G-v)\bigr)\notag\\
 &{}\lor\bigvee_{uv\in E(G)}
 \bigl(S(\Nb_G(u,v))\land S(G-uv)\bigr).
 \label{eq:strong-rec}
\end{align}
Evaluation may stop when one conjunction is true; a negative answer
requires every conjunction to be false. The argument of $S$ is always
the current graph, including the edges that remain after earlier
deletions.

\begin{proposition}\label{prop:IS-correct}
The evaluation of~\eqref{eq:strong-rec} terminates and returns true
exactly when $G\in\IS$. A true answer determines a reduction to $K_1$
and reductions of the neighbourhoods that justify its steps.
\end{proposition}
\begin{proof}
Induct on $|V(G)|+|E(G)|$. Every neighbourhood and every remaining
graph in~\eqref{eq:strong-rec} has smaller measure, and there are
finitely many alternatives. This proves termination, starting from
the two stated bases. If a conjunction is true, induction supplies
the two reductions that make the corresponding first deletion valid;
prepending that deletion proves $G\in\IS$. Conversely, a strong
reduction of a non-singleton graph has a first vertex or edge deletion.
Its neighbourhood and remainder are in $\IS$, so induction makes
the corresponding conjunction true. Recording a successful choice
at each occurrence gives the asserted reductions.
\end{proof}

For vertex reductions, all intermediate graphs are induced subgraphs
of a fixed graph $G$. For $U\subseteq V(G)$ and $v\in U$, put
$N_U(v)=U\cap V(\Nb_G(v))$. Define truth values $W(U)$ and $D_k(U)$
to be false on the empty set and true on singletons. For nonempty $U$,
$D_{-1}(U)$ means that $G[U]$ is a cone. On sets of size at least two,
use
\begin{align}
 W(U)&=\bigvee_{v\in U}
     \bigl(W(N_U(v))\land W(U\setminus\{v\})\bigr),
     \label{eq:ivs-rec}\\
 D_k(U)&=\bigvee_{v\in U}
     \bigl(D_{k-1}(N_U(v))\land D_k(U\setminus\{v\})\bigr),
     \qquad k\ge0.\label{eq:kdis-rec}
\end{align}

\begin{proposition}\label{prop:vertex-correct}
The recurrences~\eqref{eq:ivs-rec} and~\eqref{eq:kdis-rec} decide
whether $G[U]$ belongs to $\ISV$ and to $k$-Dis, respectively. For a
fixed largest level, evaluation in increasing order of $|U|$
terminates.
\end{proposition}
\begin{proof}
Each set on the right is a proper subset of $U$. Induct on its
cardinality, evaluating the cone test directly. A successful
conjunction supplies a first admissible vertex and a reduction of
the remainder. In the opposite direction, the first step of any
reduction supplies one of the conjunctions. The neighbourhood
condition is at the same recursive class for $\ISV$ and at level
$k-1$ for $k$-Dis, exactly as in their definitions. This proves both
equivalences and termination.
\end{proof}

At level zero there is also a simpler procedure: delete any dominated
vertex until no such vertex remains. Here $v$ is dominated by $w\ne v$
when $N_G[v]\subseteq N_G[w]$, where $N_G[v]=\{v\}\cup V(\Nb_G(v))$.
This is equivalent to $\Nb_G(v)$ being a cone with apex $w$.
The procedure succeeds precisely on the $0$-dismantlable graphs
\cite[Section~4.1]{FieuxJouve2020}. To see why the order does not
matter, consider two available deletions. Unless the vertices dominate
each other, each remains dominated after the other is deleted:
if a dominator is removed, use its dominator and the transitivity of
closed-neighbourhood inclusion. Mutual domination gives identical
closed neighbourhoods, so the two remainders are isomorphic.
Induction on the number of vertices now shows that an arbitrary
available deletion preserves dismantlability. Reaching $K_1$ is itself
a dismantling, proving the converse. This procedure gives the
$0$-Dis column of the table.
No analogous order-independence holds for~\eqref{eq:strong-rec}:
the extensions in Section~\ref{sec:examples} admit a deletion that
leaves a non-strong graph.

\subsection{Positive tests for I-contractibility}

Strong reductions also supply positive tests for the larger class
$\I$. Suppose
\begin{equation}\label{eq:I-path}
 G=G_0,\qquad G_j=G_{j-1}+u_jv_j\quad(1\le j\le t),
 \qquad G_t\in\IS,
\end{equation}
and every added edge has common neighbourhood in $\IS$ at its step.
Then $G\in\I$. Indeed, start from $G_t\in\IS\subseteq\I$ and delete
the added edges in reverse order. Adding or deleting the edge between
two vertices leaves their common neighbourhood unchanged, so each
reverse step is an I-transformation.

For a fixed bound $t$, enumerate all admissible edge additions from
$G$, one step at a time, and test each reached graph by
Proposition~\ref{prop:IS-correct}. This finds a path of the form
\eqref{eq:I-path} of length at most $t$ whenever one exists: induction
on the path length shows that every possible successor is considered.
The search is finite, since the vertex set is fixed and edges are
only added. Its failure excludes only the specified paths, not
membership in $\I$. The negative decisions needed below instead use
nonzero integral homology. Together these two tests will settle all
graphs left outside $\IS$ in the universes of the table.

\subsection{Representatives up to isomorphism}\label{sec:representatives}

The generation uses canonical vertex augmentation
\cite[Section~4]{McKay1998}. For each graph $H$ of order at least
two, choose a vertex orbit $O(H)$ under $\operatorname{Aut}(H)$,
invariantly under isomorphism. This choice can be made from a
canonical labelled representative. Begin with $K_1$. For every
retained graph $G$, adjoin a vertex $x$ adjacent to one subset from
each orbit of subsets of $V(G)$ under $\operatorname{Aut}(G)$.
Retain the resulting graph $H$ when $x\in O(H)$. Connectedness is
required at the output order. For the $K_4$-free universe, reject
augmentations containing a $K_4$.

\begin{proposition}\label{prop:enumeration-correct}
This construction, followed by the connectedness test, generates
exactly one representative of every connected graph of the requested
order. The additional $K_4$ rejection gives exactly the connected
$K_4$-free graphs.
\end{proposition}
\begin{proof}
Induct on the order, before testing connectedness. For a
target $H$, choose $x\in O(H)$. Its parent is represented by induction
and the neighbourhood orbit is considered. The resulting copy of
$H$ is accepted because $O$ is invariant under isomorphism, proving
existence.
For two isomorphic accepted augmentations, an isomorphism followed
by an automorphism identifies the new vertices, and hence the parents
and neighbourhood orbits. These were considered once, proving
uniqueness. Restricting the outputs preserves both conclusions.
Every induced subgraph of a $K_4$-free graph is $K_4$-free, so the
additional rejection loses no predecessor of an eligible target.
\end{proof}

For any of the exact predicates above, sum its value over the generated
representatives of a fixed order. Proposition~\ref{prop:enumeration-correct}
and the decision propositions show that this sum is the cardinality
of the corresponding class, with neither missing isomorphism classes
nor repeated contributions. The column for $\I$ additionally requires
the positive and negative adjudications described in
Section~\ref{sec:table}.
The inclusions in~\eqref{eq:dismantling-inclusions} and
\eqref{eq:homology-inclusions} allow a positive decision for one
class to establish membership in every larger class of these chains;
a negative decision establishes non-membership in every smaller class.

\subsection{Reduction of the order-fourteen classification}\label{sec:reduction}

The complete census gives the class counts. To identify the acyclic
non-strong graphs behind these counts, we reduce them until no
admissible move remains.

\begin{definition}\label{def:stuck}
A graph of order at least two is \defi{stuck} if it is
$\mathbb F_2$-acyclic and has no $\IS$-admissible vertex or edge.
A graph of order at least two is \defi{almost stuck at $u$} if it is
$\mathbb F_2$-acyclic,
has no $\IS$-admissible edge, and no vertex other than $u$ is
$\IS$-admissible. The distinguished vertex $u$ need not be admissible.
\end{definition}

Every stuck graph lies outside $\IS$, since a reduction cannot start.
Two observations explain the role of minimum order. An admissible
deletion from a graph of $\I$ remains in $\I$. Moreover, if its
remainder is in $\IS$, then the original graph is in $\IS$, by
prepending that deletion to a strong reduction.

\begin{lemma}\label{lem:stuck}
Let $\mathcal H$ be closed under vertex and edge deletion, and let $n$
be the smallest order attained in $\mathcal H\cap(\I\setminus\IS)$.
Every graph in this class of order $n$ reduces by admissible edge
deletions to a stuck graph in the same class.
\end{lemma}
\begin{proof}
An admissible vertex deletion would give a smaller graph in the same
class, by the two observations, contradicting the choice of $n$.
Delete admissible edges as long as possible. The order is unchanged,
membership in $\mathcal H\cap(\I\setminus\IS)$ is preserved, and
the number of edges decreases. At termination neither kind of
deletion is possible. The terminal graph is acyclic because it is
in $\I$, so it is stuck.
\end{proof}

\begin{lemma}\label{lem:2conn}
If a stuck graph has a cut vertex $u$, it is the union of two induced
graphs $A,B$, each of order at least three, that meet only at $u$ and
are almost stuck there. Consequently, if no almost stuck graph in a
hereditary class has order at most $m$, every stuck graph in that
class of order at most $2m$ is $2$-connected.
\end{lemma}
\begin{proof}
Partition the components of $G-u$ into two nonempty collections and
include $u$ in both pieces. Their clique complexes meet at one
vertex, so reduced Mayer--Vietoris gives
\[
 \widetilde H_i(G;\mathbb F_2)\cong
 \widetilde H_i(A;\mathbb F_2)\oplus
 \widetilde H_i(B;\mathbb F_2).
\]
Both pieces are acyclic. The neighbourhood of a vertex other than
$u$, and the common neighbourhood of every edge in a piece, agree
with those in $G$. Thus no such element is admissible. A piece of
order two would have a leaf of $G$, which is admissible. Both orders
are therefore at least three. If both exceed $m$, their union has
order at least $2(m+1)-1=2m+1$, proving the last assertion.
\end{proof}

\begin{lemma}\label{lem:degrees}
Every edge of a $2$-connected stuck graph $G$ lies in at least two
triangles, and
\[
 4\le\delta(G)\le\Delta(G)\le |V(G)|-2.
\]
\end{lemma}
\begin{proof}
An edge in no triangle belongs to a cycle, by $2$-connectedness.
Its coefficient in every $2$-boundary is zero, whereas its coefficient
in that cycle is one. The cycle therefore represents nonzero homology
over $\mathbb F_2$, contradicting acyclicity.
An edge in exactly one triangle has common neighbourhood $K_1$ and
is admissible. Hence every edge lies in at least two triangles.
For $H=\Nb_G(v)$, every vertex $w$ of $H$ consequently has at least
two neighbours in $H$. If $|V(H)|\le3$, this forces $H=K_3$, making
$v$ admissible. Thus every degree is at least four. A universal
vertex would make $G$ a cone and hence a member of $\IS$, so no
degree is $|V(G)|-1$.
\end{proof}

On $K_4$-free graphs these reductions have particularly simple tests.
\begin{proposition}\label{rem:K4-local}
In a $K_4$-free graph, a vertex is $\IS$-admissible exactly when its
open neighbourhood is a nonempty tree, and an edge is $\IS$-admissible
exactly when it belongs to one triangle.
\end{proposition}
\begin{proof}
A vertex neighbourhood is triangle-free. A triangle-free graph in
$\IS$ has a contractible clique complex of dimension at most one,
so it is a tree. Conversely, a nonempty tree reduces by leaf deletions
to $K_1$. The common neighbourhood of an edge is independent and
belongs to $\IS$ exactly when it has one vertex. Its vertices
correspond to the triangles containing the edge.
\end{proof}

This criterion is independent of any small-order census. In particular,
to test stuckness on a $K_4$-free graph, check its reduced homology and
exclude every tree vertex neighbourhood and every edge in one triangle.
For almost stuckness allow at most one such vertex, including the
possibility of none.

Finally, reverse the edge deletions of Lemma~\ref{lem:stuck} to recover
the graphs arising from a stuck graph. Starting with a finite set of
stuck graphs, called \defi{seeds}, adjoin every missing edge with common neighbourhood in
$\IS$; continue from non-strong successors and discard strong ones.
Identify isomorphic graphs throughout. The retained set is the
\defi{pruned edge-addition closure}.

\begin{proposition}\label{prop:closure-correct}
This procedure terminates and retains exactly the non-strong graphs
reachable from the seeds by admissible edge additions. To recover
only $K_4$-free targets, one may also discard successors containing
a $K_4$.
\end{proposition}
\begin{proof}
There are finitely many graphs on the fixed vertex set and every
step adds an edge. All retained graphs are reachable, by construction.
If $H\in\IS$ and $H+uv$ is an admissible addition, deleting $uv$
first and then reducing $H$ proves $H+uv\in\IS$. A path to a
non-strong target therefore never passes through a strong graph.
Every predecessor on that path is expanded, so induction on path
length proves completeness. Isomorphisms transport nonedges and
their neighbourhoods, and hence all possible continuations.
Finally, adding edges cannot remove a $K_4$, so no path to a
$K_4$-free target is lost by the additional restriction.
\end{proof}
